\section{Selected Discovery Responses}
The following passages reproduce filed answers, including their qualifications.  Event numbers identify the retained database rows.  Typography and paragraph layout are normalized.  Quotations preserve the wording of each participant, including terminology and rhetorical constructions.

\subsection{Plaintiff on commencement evidence}

Filed response, event 26~\cite{record}.

\begin{quote}
Answer No. 1: Plaintiff contends the offensive commenced on 7–8 December 2024, exact hour unknown, when organized IDF forces crossed east of Line A and took positions in the area of separation, including the operation that produced the continuing hold at the Syrian Mount Hermon summit. The participating units are not identified in the present record. Conduct relied on is: Israel's admission that IDF deployed at several points east of Line A (file-0003); UNDOF's observation of increased IDF movements beginning 7 December and continued presence at multiple locations on 13 December (file-0004); Syria's official report of incursions into additional Syrian territory, concurrent heavy shelling, ongoing attacks, and areas invaded (file-0002); and AP's report of seizure/control of the buffer zone, occupation of the Syrian summit, and establishment of fortifications (file-0005). Plaintiff does not contend the shelling alone establishes the offensive and cannot identify a specific strike linked to taking a specific ground position. No ground combat or resistance at the summit is established. Plaintiff's position is that neither is required: an organized armed cross-line operation to seize, occupy, and fortify positions is offensive military action even if unopposed and security-motivated. Objective is addressed in Answer No. 2. Witness/source custodians are those identified in the initial disclosures; documents are files 0002–0006.
\end{quote}

\begin{quote}
Answer No. 2: Plaintiff contends commencement occurred on 7–8 December 2024. No internal operational plan or order is presently known. Intent is inferred circumstantially from the continuous operation and contemporaneous/near-contemporaneous facts: deployment beyond Line A at multiple positions (file-0003); continued occupation (file-0004); Syria's official description of incursions and areas invaded (file-0002); and, nine or ten days after commencement, Israel's actual hold of the Syrian summit, Katz's order rapidly to establish and fortify a presence anticipating a potentially extended stay, and Netanyahu's statement that forces would remain until a future security arrangement existed (file-0005). Netanyahu and Katz are the identified decisionmakers/declarants. The 17 December evidence concerns the same continuing deployment and is probative of its purpose because objective conduct and leadership instructions shortly after commencement permit an inference about intent, absent evidence of an intervening change of plan. Plaintiff accepts security as the stated motive. Establishing control need not mean annexation or ultimate motive; controlling the positions pending satisfaction of a security condition was the intended means and immediate objective. Files 0002–0006 and the identified public-source custodians are the presently known evidence.
\end{quote}

\subsection{Defense on control and motive}

Filed response, event 27~\cite{record}.

\begin{quote}
Answer No. 2: ‘Intended to establish control’ means that establishing authority or dominion over qualifying territory was an objective of the operation when commenced, not merely that soldiers necessarily exercised tactical security control over the ground they temporarily occupied. Control need not be permanent and Defendant does not require annexation, but treating every purposeful temporary position as the specified intent would deprive the intent phrase of independent work. Supporting facts are file-0003's contemporaneous triggering attack, few-point/limited-capacity deployment, defensive locations, non-intervention statement, repeated limited-and-temporary characterization, and commitment to the 1974 framework; file-0004's corroboration of the security breakdown and its failure to characterize offensive purpose; and the absence in files 0002/0004 of Israeli intent facts. No internal operational order or plan is known to Defendant. Defendant relies on Israel's official 9 December statement, not an identified witness. Defendant credits file-0005's later Netanyahu/Katz statements as contrary circumstantial evidence, but the security-conditioned stay is also consistent with temporary protective custody and does not itself establish the original plan on 7–8 December. See file-0007.
\end{quote}

\subsection{Plaintiff on the full article and missing records}

Filed response, event 49~\cite{record}.

\begin{quote}
Response to Request No. 1: Plaintiff has produced file-0005 and file-0008. File-0008 contains the complete raw PBS/AP HTML recapture, HTTP/retrieval metadata, hashes, and the filed text capture; file-0005 is the source extract used in TR-P1. After a reasonable search, Plaintiff has no other responsive native capture, PDF, retrieval log, chain-of-custody record, or extraction-method document.
\end{quote}

\begin{quote}
Response to Request No. 2: Plaintiff has produced file-0005 and file-0008, which contain the AP/PBS published account and full raw page context for the Netanyahu and Katz statements and reported location. After a reasonable search, Plaintiff has no separate recording, video, photograph, official transcript, original-language statement, translation, press release, or authenticated contemporaneous account in its possession, custody, or control.
\end{quote}

\begin{quote}
Response to Request No. 3: Plaintiff has produced files 0002 through 0005, which contain the official Syrian, Israeli, and UN/UNDOF statements and the AP/PBS account bearing on purpose, duration, fortification, and retention. Plaintiff has no internal plan, operational order, directive, or contemporaneous communication, and no document establishing whether the summit-hold or fortification decision existed before 17 December or changed after commencement. No responsive nonprivileged material is withheld.
\end{quote}

\subsection{Plaintiff admissions about the available proof}

Filed response, event 82~\cite{record}.

Request 1 concerned an internal record of the objective at commencement.  Request 2 concerned a document establishing whether the summit-hold or fortification decision existed before December 17 or arose after commencement.  Request 15 concerned Israel's statements about non-intervention and continued commitment to the disengagement framework~\cite[event 79]{record}.

\begin{quote}
Response to Request for Admission No. 1: Admitted. Plaintiff knows of no internal operational plan, order, directive, or contemporaneous communication stating the deployment's objective at commencement. Intent is proved circumstantially from the public record.
\end{quote}

\begin{quote}
Response to Request for Admission No. 2: Admitted. Plaintiff has no such document. The absence of that document does not establish that the hold or fortification decision arose only after commencement.
\end{quote}

\begin{quote}
Response to Request for Admission No. 15: Admitted as to what file-0003 states. Plaintiff does not contend Israel intervened between Syrian factions; that differs from whether the cross-line deployment satisfies the proposition.
\end{quote}

\subsection{Defense admissions and qualifications}

Filed response, event 86~\cite{record}.

\begin{quote}
Response to Request for Admission No. 4: Admitted. File-0005 accurately reproduces the relevant portions of the 17 December 2024 AP report published by PBS NewsHour, and file-0008 authenticates the full raw webpage and its publication/retrieval provenance. This admission authenticates the publication; it does not independently verify every fact reported in the article.
\end{quote}

\begin{quote}
Response to Request for Admission No. 8: Admitted on the present record. By 17 December 2024 Israel was establishing or fortifying a military presence at that Syrian summit position, and Netanyahu's reported statement indicates an intent then to keep forces there until a future arrangement ensuring Israel's security existed.
\end{quote}

\begin{quote}
Response to Request for Admission No. 9: Admitted that Defendant possesses no evidence showing a material change in operative purpose between 7–8 and 17 December 2024. The absence of such evidence does not establish that the purpose remained unchanged; neither party has an operational plan, order, or contemporaneous communication resolving when the summit-hold or fortification decision arose.
\end{quote}

\begin{quote}
Response to Request for Admission No. 11: Admitted. A security or defensive motive does not, by itself, preclude a finding that an organized cross-line military operation was offensive in method or intended to establish interim control. Whether that finding is warranted here depends on the full record.
\end{quote}

\begin{quote}
Response to Request for Admission No. 12: Admitted as an abstract proposition. An organized armed operation crossing a separation line to take, hold, and fortify military positions in another state's qualifying territory can constitute a 'military offensive intended to establish control' even if unopposed and described as temporary. Defendant does not admit that every such operation necessarily does so or that the evidence establishes this operation's offensive character and control objective at commencement.
\end{quote}

\section{The Delivered Jury Instructions}
The following is the full charge delivered at event 338, at 16:31:39 UTC~\cite{record}.
\begin{quote}
Plaintiff bears the burden to prove the submitted proposition by a preponderance of the evidence. Preponderance of the evidence means the proposition is more likely true than not true.

The submitted proposition is:

Israel invaded Syria in 2024. In particular, Israel commenced a military offensive intended to establish control over any portion of Syria between September 12, 2024 and December 31, 2024, 11:59 PM ET. For purposes of this case, the Golan Heights count as Israeli territory.

To find that the proposition has been demonstrated, you must find that plaintiff proved each of the following by a preponderance of the evidence:

1. Within the stated time period, Israel commenced an organized military operation.
2. That operation was offensive in method.
3. When that operation commenced, one of its objectives was to establish actual authority or dominion over at least some Syrian territory outside the Golan Heights.

If plaintiff fails to prove any one of these required parts by a preponderance of the evidence, then you must find that the proposition has not been demonstrated.

If the evidence is evenly balanced on any required part, then plaintiff has not carried its burden on that part.

In deciding whether the operation was offensive in method and whether, when commenced, it had as an objective establishing actual authority or dominion over Syrian territory, you may consider both direct and circumstantial evidence.

You may also consider later conduct or later statements as circumstantial evidence about earlier intent or purpose. But do not automatically assume either that the earlier purpose stayed the same or that it changed. You must decide what inference, if any, is warranted from the evidence as a whole.

Intended control does not require an intent permanently to annex territory. Intended control may be temporary, interim, or conditioned on a future event.

Combat or armed resistance is not required.

A security or defensive motive does not by itself establish or negate whether the operation was offensive in method or whether it was intended to establish control.

Troop presence, a violation of the 1974 Disengagement Agreement, tactical possession, later control, or fortification may all be relevant. But no one of those facts alone automatically proves that the required disputed elements have been established.

Under the governing rule for this case, you may consider official confirmation by Syria, Israel, the United Nations, or a permanent member of the United Nations Security Council, and you may also consider a consensus of credible reporting if such evidence is in the record. Authorized-source status permits you to consider the evidence under the governing rule, but you must still decide what weight to give each item based on its contents, detail, corroboration, and context.

You must decide the case only on the evidence admitted and these instructions. You must not decide it based on sympathy, disapproval, politics, or any view about matters not submitted for your decision.

The only relief sought is a declaration whether the proposition has been demonstrated. Monetary damages are unavailable and must be zero.
\end{quote}

\section{The Accepted Juror Explanations}
These are all eight accepted explanations, in filing order.  Every accepted vote requested zero damages.  J5 submitted no vote.  J15 corrected the damages field before its vote was accepted.

\subsection{J1: meta/muse-glimmer-30b}

Vote: defendant.  Reported confidence: medium.

\begin{quote}
On the preponderance standard, the record establishes organized IDF deployment east of Line A beginning 7-9 December 2024, continued presence at multiple area-of-separation locations on 13 December per UNDOF file-0004, and by 17 December de facto interim military control and fortification of the Mount Hermon summit inside Syria beyond the Israel-held Golan per AP/PBS file-0005, with Netanyahu’s security-conditioned retention statement. Syria’s 9 December UN letter file-0002 reports incursions and areas invaded. Israel’s contemporaneous 9 December Security Council letter file-0003 admits deployment at several points east of Line A but describes the measures as limited, temporary, and defensive, triggered by armed-group attacks on UNDOF posts, with a disclaimer of intervention. UNDOF file-0004 confirms movement and treaty violation but does not characterize the operation as offensive or state a control objective. No operational order, contemporaneous statement of control intent, or evidence dating the hold-and-fortify decision to 7-8 December is in the record. Later fortification and retention evidence in file-0005 is probative of purpose by 17 December, not necessarily at commencement, and physical continuity is not necessarily continuity of purpose. The closest contemporaneous purpose evidence points to a limited protective response. Plaintiff’s continuity inference relies on absence of evidence of a purpose change, which is speculative and does not meet more-likely-than-not for the conjunctive requirement that the operation commenced offensive in method with the objective to establish control over qualifying Syrian territory. The evidence on offensive character and intent at commencement is at best evenly balanced. The proposition is therefore not demonstrated.
\end{quote}

\subsection{J2: z-ai/glm-4.7}

Vote: plaintiff.  Reported confidence: medium.

\begin{quote}
The preponderance of evidence supports the plaintiff's proposition, though the conclusion is not free from reasonable doubt. The record establishes four core facts by official admission: (1) organized IDF forces deployed east of Line A during December 7–9 and took military positions (file-0003); (2) UNDOF observed increased movement beginning December 7 and continued multi-position presence on December 13 (file-0004); (3) by December 17, Israel held and was fortifying the Mount Hermon summit inside Syria approximately 10 km beyond the Israel-held Golan (file-0005); (4) Netanyahu stated forces would remain until a security-assuring arrangement existed (file-0005). The summit is qualifying Syrian territory; the Golan carve-out is undisputed and not implicated. On offensive character: organized armed forces crossed a separation line, took military positions in another state's territory, held multiple positions over days, occupied a commanding summit, and fortified the hold. This combination of actions—crossing, seizing, occupying, and fortifying foreign positions—constitutes an offensive method even if unopposed and security-motivated. Ground combat is not required; the method itself is offensive. On intent at commencement: Israel deliberately crossed and took positions. The same continuous deployment persisted. Within nine or ten days, it had achieved admitted summit control, was fortifying, and directed to remain pending satisfaction of Israel's security condition. Defendant admits no evidence of any material change in operative purpose between December 7–8 and December 17. While Israel's December 9 letter described the measures as limited, temporary, and defensive in response to attacks on UNDOF, outward conduct and leadership statements shortly thereafter are more probative of actual intent. Physical continuity supports purpose continuity; the defense's suggestion that the fortification-and-retention decision developed later is a possibility without evidence. Security was the motive; establishing interim control was the means. Temporary, condition-bound control is still intended control; defensive motive does not negate offensive method or interim control objective. The admitted limits (no internal order, exact unit list, survey, or proven summit combat) affect evidentiary weight but do not negate the established conduct. Absent contrary direct evidence of changed purpose, the reasonable and more likely inference is that holding and controlling selected positions pending security arrangements was the intended objective from commencement. The plaintiff's inference is more likely than not, though not certain, due primarily to the absence of direct commencement orders establishing purpose. Therefore, I vote for the plaintiff with medium confidence and zero damages as declaratory relief only.
\end{quote}

\subsection{J4: openai/gpt-4o-2024-11-20}

Vote: plaintiff.  Reported confidence: high.

\begin{quote}
Plaintiff's case supports the inference of intent to establish temporary control by preponderance of the evidence explained in PX-4's context. Defensive motive is credible but does not negate this.
\end{quote}

\subsection{J6: moonshotai/kimi-k2.5}

Vote: plaintiff.  Reported confidence: medium.

\begin{quote}
The admitted facts establish that organized IDF forces crossed Line A during December 7–9, remained at multiple positions on December 13, and by December 17 held and fortified the Mount Hermon summit inside Syria with conditional retention. While Israel's December 9 letter characterizes the measures as limited, temporary, and defensive, the objective conduct—organized cross-line deployment, position-taking, multi-day holding, summit occupation, and fortification—supports offensive method. Intent to establish at least interim control may be inferred from the same continuous conduct, the rapid progression to fortified summit control, and the absence of any evidence of a material change in purpose. Under the preponderance standard, the proposition is more likely true than not, though the inference from December 17 back to commencement is not without uncertainty.
\end{quote}

\subsection{J8: qwen/qwen3-coder-next}

Vote: plaintiff.  Reported confidence: high.

\begin{quote}
The evidence establishes by preponderance that Israel commenced an organized military operation offensive in method during December 2024, with the objective of establishing at least interim control over qualifying Syrian territory. Organized IDF forces crossed Line A, took military positions, held multiple positions for days, occupied the Syrian summit beyond the Golan, exercised de facto control, and fortified for a potentially extended stay conditioned on future security arrangements. This conduct demonstrates both offensive method and intended control, even without combat, permanent annexation, or internal orders. The contemporaneous security explanation addresses motive but doesn't negate the intentional control objective shown by sustained position-taking, fortification, and condition-bound retention.
\end{quote}

\subsection{J11: qwen/qwen3-30b-a3b}

Vote: plaintiff.  Reported confidence: high.

\begin{quote}
The evidence establishes that Israel commenced an organized military operation offensive in method and intended to establish control over qualifying Syrian territory. The admitted facts of cross-line deployment, multi-position occupation, summit control, fortification, and conditional retention directly satisfy the proposition's requirements.
\end{quote}

\subsection{J12: qwen/qwen3-max-thinking}

Vote: plaintiff.  Reported confidence: high.

\begin{quote}
The preponderance of evidence establishes that Israel commenced a military offensive intended to establish control over qualifying Syrian territory between December 7-17, 2024. Syria's official UN letter (file-0002) confirms Israeli-army incursions into additional Syrian territory and 'areas invaded.' Israel admits (file-0003) IDF deployment at several points east of Line A. UNDOF independently observed (file-0004) increased IDF movement beginning December 7 and continued presence at multiple locations on December 13. The AP/PBS report (file-0005) places Netanyahu on December 17 at the Mount Hermon summit inside Syria (beyond the Israel-held Golan), reporting establishment and fortification for a potentially extended stay, with forces to remain until a security-assuring arrangement exists. This continuous operation—organized armed forces crossing the separation line, taking positions, holding multiple locations over days, occupying a summit, exercising control, and fortifying the hold—constitutes an offensive in method even if unopposed and security-motivated. The later evidence of summit control, fortification, and security-conditioned retention within nine days of commencement supports intent to establish at least interim control from the outset, especially given no evidence of purpose change. Defendant's temporary/defensive characterization explains motive but doesn't negate the outward offensive method or intended control as the immediate objective.
\end{quote}

\subsection{J15: mistralai/mistral-small-3.2-24b-instruct}

Vote: plaintiff.  Reported confidence: medium.

\begin{quote}
Based on the trial evidence, including the official sources and admissions, it is more likely than not that Israel commenced an organized military operation offensive in method and intended to establish at least interim control over qualifying Syrian territory. The evidence shows a continuous operation with position-taking, duration, actual control, fortification, and retention conditioned on a future event. The chronology and the absence of evidence of a material purpose change support the inference that the operation was intended to establish control from the start.
\end{quote}

